% =============================================================================
% 066 / 068
% Status: raw
% =============================================================================
% Notes:
%
% Source question:
% 其實，早就有臺語(閩南語)或其他各東亞大陸語言的使用者，
% 零星地整理他們母語的語法。
% 
% 但是，我引用推友 並明 所說的：
% 
% 『用”漢語語法”的體系來整理語法，
% 而不是原生語法整理語法。
% 其實是打著”整理”旗號做藍青化的工整吧。』
% 
% 整理語法易，避開官話思維整理語法難，
% 這是一個很大的陷阱？
% =============================================================================

\chapter[其實，早就有臺語(閩南語)或其他各東亞大陸語言的使用者， 零星地整理他們母語的語法。 但是，我引用推友 並明 所說的： 『用”漢語語法”的體系來整理語法， 而不是原生語法整理語法。 其實是打著”整理”旗號做藍青化的工整吧。』 整理語法易，避開官話思維整理語法難， 這是一個很大的陷阱？]{其實，早就有臺語(閩南語)或其他各東亞大陸語言的使用者， \\[0.35em] 零星地整理他們母語的語法。 \\[0.35em] 但是，我引用推友 並明 所說的： \\[0.35em] 『用”漢語語法”的體系來整理語法， \\[0.35em] 而不是原生語法整理語法。 \\[0.35em] 其實是打著”整理”旗號做藍青化的工整吧。』 \\[0.35em] 整理語法易，避開官話思維整理語法難， \\[0.35em] 這是一個很大的陷阱？}
\qaanswer{
整理語法不易，現代漢語普通話語法和東亞大陸多數地區的原生語言，包括官話區內的各亞方言，本身就存在差異，只是現代漢語普通話與這些方言、語言之間有些僅僅是量的差異，有些則是質的不同。以吳語為例，她是一種與漢語族Mandarin語支 下的各語言，語言思維(語灋)差異非常大的獨立語言。她雖然收到漢語族Mandarin語支的強烈影響，她的一些方言，如上海話，已經呈現和Mandarin語灋越來越趨同的現象，但她最核心的語言思維還是普遍存在於吳語區多數地區的中高齡使用者腦海中的。
}

